\documentclass[a4paper,12pt,openany]{book}
\input{../../../preambule/preambule_packages}
\input{../../../preambule/preambule_mise_en_page}
\input{../../../preambule/preambule_math}
\input{../../../preambule/preambule_environnements}
% en-tête

\pagestyle{otherpage}
\fancyhead[R]{1NSI}
\fancyhead[L]{Nom - Prénom}
\begin{document}

\begin{center}\begin{large}
DS02 - Structures FOR / IF
\end{large}\end{center}

\exo

\begin{minipage}{9cm}
	\begin{lstlisting}[numbers=left]
def f(a,b):
		if a < b:
				return a
		else:
				return b
\end{lstlisting}
\end{minipage}
\begin{minipage}{9cm}
\begin{liste}
	\item 	Que renvoie la commande \verb+f(3,7)+ ? \dotfill
	\item Dans quel ordre sont exécutées les lignes du programme lors de l'appel \verb+f(3,7)+ ? 
	
	\dotfill
	\end{liste}
\end{minipage}	

	
	
	\bigskip

\exo

\begin{minipage}{9cm}
	\begin{lstlisting}[numbers=left]
def g(x):
		for i in range(4):
				x = x + i
		return x
\end{lstlisting}

\begin{liste}
	\item 	Que renvoie la commande \verb+g(2)+ ? \dotfill
	
	\dotfill
	\end{liste}\end{minipage}
\begin{minipage}{9cm}

\begin{liste}
	\item Dans le tableau suivant, donner l'évolution des différentes variables lors de l'appel \verb+g(2)+
	
\renewcommand{\arraystretch}{1.2}	
\begin{tabular}{|p{2cm}|p{2cm}|p{2cm}|}
\hline
ligne & \verb+i+ & \verb+x+  \\ \hline
1 & XXX & 2\\ \hline
2 & 0 & \\ \hline
& & \\ \hline
& & \\ \hline
& & \\ \hline
& & \\ \hline
& & \\ \hline
& & \\ \hline
& & \\ \hline
& & \\ \hline
& & \\ \hline

\end{tabular}	
	
\end{liste}\end{minipage}	






	
	\bigskip
	
\exo 


\begin{minipage}{9cm}
\begin{lstlisting}[numbers=left]
def h(mot):
		# mot: chaine de caractères ne contenant que des chiffres. Par exemple '163892' 
		a = 0
		for valeur in mot:
				if int(valeur) > a:
						a = int(valeur)
		return a
		
\end{lstlisting}

\begin{liste}
	\item 	Que renvoie la commande \verb+h('2634')+ ? \dotfill
	
	\dotfill
	
	\item Quel est le type de la variable \verb+a+ ? \dotfill
	
	\item Quel est le type de la variable \verb+ valeur + ? \dotfill
	 
	\end{liste}
\end{minipage}
\begin{minipage}{9cm}
\begin{liste}
	\item Dans le tableau suivant, donner l'évolution des différentes variables lors de l'appel \verb+h('283')+
	
\renewcommand{\arraystretch}{1.2}	
\begin{tabular}{|p{1.5cm}|p{1cm}|p{2cm}|p{2cm}|}
\hline
ligne & \verb+a+ & valeur & IF  \\ \hline
3 & 0 & XXX&XXX\\ \hline
4 &  & &XXX\\ \hline
5 & & &\\ \hline
 & & &\\ \hline
 & & &\\ \hline
 & & &\\ \hline
 & & &\\ \hline
 & & &\\ \hline
 & & &\\ \hline
 & & &\\ \hline
 & & &\\ \hline
 & & &\\ \hline
 & & &\\ \hline


\end{tabular}	
	
\end{liste}
\end{minipage}	




\newpage

\exo 

\begin{minipage}{10cm}
\begin{lstlisting}[numbers=left]
def fonction(n):
		# n entier strictement positif 
		s = 0
		for i in range(6):
				if i >= n:
						s = s + i
		return s
		
\end{lstlisting}
\end{minipage}
\begin{minipage}{8cm}
\begin{liste}
	\item 	Que renvoie la commande \verb+fonction(4)+ ? \dotfill
	
	\dotfill
	
	 	\item Dans le tableau suivant, donner l'évolution des différentes variables lors de l'appel \verb+fonction(4)+

	\end{liste}
\end{minipage}
	




	
\begin{center}
\renewcommand{\arraystretch}{1.2}	
\begin{tabular}{|p{2cm}|p{2cm}|p{2cm}|p{2cm}|}
\hline
ligne & \verb+s+ & \verb+i+ & \verb+IF i >= n+    \\ \hline
3 & 0 & XXX&XXX\\ \hline
4 &  &0 &XXX\\ \hline
5 & & &\\ \hline
 & & &\\ \hline
 & & &\\ \hline
 & & &\\ \hline
 & & &\\ \hline
 & & &\\ \hline
 & & &\\ \hline
 & & &\\ \hline
 & & &\\ \hline
 & & &\\ \hline
 & & &\\ \hline
 & & &\\ \hline
 & & &\\ \hline
 & & &\\ \hline
 & & &\\ \hline


\end{tabular}	
\end{center}
	
	
	
	
	
\exo

\begin{minipage}{8.5cm}
Écrire la fonction \verb+compte_a(mot)+ qui renvoie le nombre de \verb+a+ dans \verb+mot+. Par exemple \verb+compte_a('abracadabra')+ renverra 5.

\begin{lstlisting}[numbers=left]
def compte_a(mot):












#
		
\end{lstlisting}
\end{minipage}
\begin{minipage}{1cm}
||
\end{minipage}
\begin{minipage}{8.5cm}
Écrire la fonction \verb+somme_n(n)+ qui renvoie la somme des \verb+n+ premiers entiers naturels. Par exemple \verb+somme_n(6)+ renverra la somme de \verb-1+2+3+4+5+6- soit la valeur 21.

\begin{lstlisting}[numbers=left]
def somme_n(n):











#
		
\end{lstlisting}
\end{minipage}
\newpage


\end{document}
